% =====================================================================
%  Оформление по ГОСТ 7.32-2017
%  Структурные элементы: разделы, рисунки, таблицы, формулы, списки
% =====================================================================

% Размер основного шрифта 14 pt с межстрочным интервалом 1.5 устанавливается
% в classfile (extarticle/extreport, базовый размер 14pt). Здесь — стилевые
% настройки заголовков и нумерации.

% Заголовки по ГОСТ 7.32-2017 (п. 6.2):
%   - заголовки разделов и подразделов с абзацного отступа,
%     с прописной буквы, без точки в конце, не подчёркиваются;
%   - расстояние между заголовком и текстом — не менее одного интервала;
%   - заголовки разделов набираются жирным шрифтом, размер 14 pt.

\titleformat{\section}
    {\normalfont\fontsize{14}{17}\bfseries}
    {\thesection}{1em}{}
\titlespacing*{\section}{\parindent}{18pt}{12pt}

\titleformat{\subsection}
    {\normalfont\fontsize{14}{17}\bfseries}
    {\thesubsection}{1em}{}
\titlespacing*{\subsection}{\parindent}{14pt}{10pt}

\titleformat{\subsubsection}
    {\normalfont\fontsize{14}{17}\bfseries}
    {\thesubsubsection}{1em}{}
\titlespacing*{\subsubsection}{\parindent}{12pt}{8pt}

% Нумерация рисунков и таблиц — сквозная в пределах раздела
\renewcommand{\thefigure}{\arabic{section}.\arabic{figure}}
\renewcommand{\thetable}{\arabic{section}.\arabic{table}}
\renewcommand{\theequation}{\arabic{section}.\arabic{equation}}

% Подписи рисунков и таблиц по ГОСТ:
%   - «Рисунок N — Название» (тире), снизу, по центру;
%   - «Таблица N — Название», сверху, слева.
\captionsetup[figure]{
    labelsep=endash,
    name=Рисунок,
    justification=centering,
    singlelinecheck=true,
    font=normalsize
}
\captionsetup[table]{
    labelsep=endash,
    name=Таблица,
    justification=raggedright,
    singlelinecheck=false,
    font=normalsize
}

% Колонтитулы — простые, номер страницы по центру нижнего поля
\usepackage{fancyhdr}
\pagestyle{fancy}
\fancyhf{}
\fancyfoot[C]{\thepage}
\renewcommand{\headrulewidth}{0pt}
\renewcommand{\footrulewidth}{0pt}

% Оглавление: «СОДЕРЖАНИЕ» по центру, прописными
\renewcommand{\contentsname}{\hfill СОДЕРЖАНИЕ\hfill}

% Список литературы: «СПИСОК ИСПОЛЬЗОВАННЫХ ИСТОЧНИКОВ»
\defbibheading{bibliography}[Список использованных источников]{%
    \section*{#1}\addcontentsline{toc}{section}{#1}%
}
